\documentclass[12pt,a4paper]{article}
\usepackage[utf8]{inputenc}
\usepackage[spanish,german]{babel}
\usepackage[T1]{fontenc}
\usepackage{geometry}
\geometry{a4paper, margin=2.5cm}
\usepackage{titlesec}
\usepackage{enumitem}
\usepackage{multicol}
\usepackage{csquotes}

\titleformat{\section}
{\Large\bfseries}
{}
{0em}
{}[\titlerule]

\titleformat{\subsection}
{\large\bfseries}
{}
{0em}
{}

\setlength{\parindent}{0pt}
\setlength{\parskip}{0.5em}

\begin{document}

\title{\textbf{DIALOG: Hablar de Salud y Ejercicio}}
\author{Spanisch A1/A2}
\date{}
\maketitle

\textbf{Tema:} Hablar de salud y ejercicio \\
\textbf{Nivel:} A2 \\
\textbf{Complejidad:} Mittel \\
\textbf{Palabras:} 856

\vspace{1em}

\textit{Nota: Se recomienda revisar primero el vocabulario antes de leer el diálogo.}

\vspace{1em}

\section*{Diálogo}

\textbf{Andrés:} ¡Hola, Sofía! ¿Cómo estás? Te veo muy bien últimamente.

\textbf{Sofía:} ¡Hola, Andrés! Gracias, me siento muy bien. He estado haciendo mucho ejercicio.

\textbf{Andrés:} ¡Qué bueno! ¿Qué tipo de ejercicio estás haciendo?

\textbf{Sofía:} Me levanto temprano todas las mañanas y voy corriendo al parque. Corriendo me siento más enérgica durante todo el día.

\textbf{Andrés:} ¿A qué hora te levantas?

\textbf{Sofía:} Me despierto a las seis de la mañana, me levanto inmediatamente y me ducho. Después me visto y salgo corriendo.

\textbf{Andrés:} ¡Eso es muy temprano! Yo me despierto a las siete, pero me quedo en la cama un rato más. Me cuesta mucho levantarme temprano.

\textbf{Sofía:} Al principio también me costaba, pero ahora me acostumbré. Haciendo ejercicio regularmente, duermo mejor y me despierto más descansada.

\textbf{Andrés:} ¿Cuánto tiempo corres?

\textbf{Sofía:} Corro durante cuarenta minutos. Primero camino un poco para calentarme, luego empiezo a correr más rápido. Terminando el ejercicio, me estiro un poco.

\textbf{Andrés:} ¿Y qué más haces para mantenerte saludable?

\textbf{Sofía:} Además de correr, voy nadando dos veces por semana. Nadando trabajo todos los músculos del cuerpo. También me gusta hacer yoga. Haciendo yoga, me relajo y reduzco el estrés.

\textbf{Andrés:} Eso suena muy completo. ¿Y qué comes? ¿Te alimentas bien?

\textbf{Sofía:} Sí, me alimento muy bien. Desayuno frutas y avena. Almorzando verduras y proteínas, me siento satisfecha. Cenando temprano y ligero, duermo mejor.

\textbf{Andrés:} Yo también quiero estar más saludable. ¿Puedes darme algunos consejos?

\textbf{Sofía:} Claro, con mucho gusto. Primero, puedes empezar caminando todos los días. Caminando media hora al día, ya estás haciendo algo bueno para tu salud.

\textbf{Andrés:} ¿Y qué más puedo hacer?

\textbf{Sofía:} Puedes beber mucha agua durante el día. Bebiendo agua, tu cuerpo funciona mejor. También puedes comer más frutas y verduras. Comiendo alimentos saludables, te sentirás mejor.

\textbf{Andrés:} ¿Y si no tengo tiempo para hacer ejercicio?

\textbf{Sofía:} Puedes hacer ejercicio en casa. Haciendo ejercicios simples en casa, también te mantienes activo. Por ejemplo, puedes hacer flexiones o sentadillas mientras estás viendo televisión.

\textbf{Andrés:} Eso es una buena idea. ¿Y cómo puedo dormir mejor?

\textbf{Sofía:} Puedes acostarte a la misma hora todas las noches. Acostándote temprano, descansas más. También puedes evitar el celular antes de dormir. Evitando las pantallas, duermes mejor.

\textbf{Andrés:} ¿Y qué haces cuando te sientes estresada?

\textbf{Sofía:} Cuando me siento estresada, respiro profundamente. Respirando profundamente, me calmo. También me relajo escuchando música tranquila. Escuchando música, me siento mejor.

\textbf{Andrés:} Eso me parece muy útil. ¿Y cómo te mantienes motivada?

\textbf{Sofía:} Me pongo metas pequeñas. Poniéndome metas alcanzables, me siento motivada. También me premio cuando logro algo. Premiándome, me mantengo entusiasmada.

\textbf{Andrés:} ¿Y qué pasa si un día no tengo ganas de hacer ejercicio?

\textbf{Sofía:} Eso es normal. Puedes hacer algo más suave ese día. Haciendo algo ligero, no pierdes el hábito. Lo importante es no dejar de moverse completamente.

\textbf{Andrés:} Entiendo. ¿Y cuánto tiempo tardaste en ver resultados?

\textbf{Sofía:} Después de un mes haciendo ejercicio regularmente, empecé a sentirme más fuerte. Después de tres meses, noté cambios más grandes. Ahora, después de seis meses, me siento completamente diferente.

\textbf{Andrés:} Eso es muy inspirador. Creo que voy a empezar mañana mismo.

\textbf{Sofía:} ¡Excelente! Recuerda: empezando poco a poco, puedes lograr grandes cosas. Si necesitas ayuda, puedes contar conmigo.

\textbf{Andrés:} Muchas gracias, Sofía. Tu ayuda significa mucho para mí.

\textbf{Sofía:} De nada, Andrés. Estoy segura de que puedes hacerlo. ¡Mucha suerte!

\newpage

\section*{Vocabulario}

\begin{multicols}{2}
\begin{itemize}[leftmargin=*]
    \item \textbf{salud} - Gesundheit
    \item \textbf{ejercicio} - Sport, Übung
    \item \textbf{últimamente} - letztens, in letzter Zeit
    \item \textbf{enérgico/a} - energisch, energiegeladen
    \item \textbf{despertarse} - aufwachen
    \item \textbf{levantarse} - aufstehen
    \item \textbf{ducharse} - duschen
    \item \textbf{vestirse} - sich anziehen
    \item \textbf{acostumbrarse} - sich gewöhnen
    \item \textbf{calentarse} - sich aufwärmen
    \item \textbf{estirarse} - sich dehnen, strecken
    \item \textbf{mantener} - aufrechterhalten, halten
    \item \textbf{mantenerte saludable} - gesund bleiben
    \item \textbf{nadar} - schwimmen
    \item \textbf{músculos} - Muskeln
    \item \textbf{yoga} - Yoga
    \item \textbf{relajarse} - sich entspannen
    \item \textbf{estrés} - Stress
    \item \textbf{alimentarse} - sich ernähren
    \item \textbf{avena} - Haferflocken
    \item \textbf{proteínas} - Proteine, Eiweiß
    \item \textbf{satisfecho/a} - zufrieden, satt
    \item \textbf{ligero/a} - leicht
    \item \textbf{consejos} - Ratschläge
    \item \textbf{flexiones} - Liegestütze
    \item \textbf{sentadillas} - Kniebeugen
    \item \textbf{activo/a} - aktiv
    \item \textbf{evitar} - vermeiden
    \item \textbf{pantallas} - Bildschirme
    \item \textbf{respirar} - atmen
    \item \textbf{profundamente} - tief
    \item \textbf{calmarse} - sich beruhigen
    \item \textbf{ponerse metas} - sich Ziele setzen
    \item \textbf{alcanzable} - erreichbar
    \item \textbf{premiarse} - sich belohnen
    \item \textbf{entusiasmado/a} - begeistert
    \item \textbf{hábito} - Gewohnheit
    \item \textbf{inspirador/a} - inspirierend
    \item \textbf{lograr} - erreichen, schaffen
    \item \textbf{contar con} - zählen auf, sich verlassen auf
\end{itemize}
\end{multicols}

\newpage

\section*{Preguntas de Comprensión}

\begin{enumerate}[leftmargin=*]
    \item ¿Cómo se siente Sofía últimamente?
    \item ¿Qué hace Sofía todas las mañanas?
    \item ¿A qué hora se despierta Sofía?
    \item ¿Qué hace Sofía después de despertarse?
    \item ¿Cuánto tiempo corre Sofía?
    \item ¿Qué hace Sofía antes de empezar a correr?
    \item ¿Qué otros ejercicios hace Sofía además de correr?
    \item ¿Por qué le gusta a Sofía hacer yoga?
    \item ¿Qué come Sofía para desayunar?
    \item ¿Qué consejo le da Sofía a Andrés para empezar?
    \item ¿Qué puede hacer Andrés si no tiene tiempo para hacer ejercicio?
    \item ¿Qué consejo le da Sofía a Andrés para dormir mejor?
    \item ¿Qué hace Sofía cuando se siente estresada?
    \item ¿Cómo se mantiene Sofía motivada?
    \item ¿Qué puede hacer Andrés si un día no tiene ganas de hacer ejercicio?
    \item ¿Cuánto tiempo tardó Sofía en ver resultados?
    \item ¿Qué cambios notó Sofía después de tres meses?
    \item ¿Qué le dice Sofía a Andrés sobre empezar?
    \item ¿Por qué Andrés quiere estar más saludable?
    \item ¿Qué significa para Andrés la ayuda de Sofía?
\end{enumerate}

\newpage

\section*{Verbo Irregular: PODER (können)}

\textit{Nota: \enquote{vosotros/as} se usa solo en España, no en español sudamericano. Se incluye aquí para completitud.}

\begin{multicols}{2}

\subsection*{Presente}
\begin{itemize}[leftmargin=*]
    \item yo - puedo
    \item tú - puedes
    \item él/ella/usted - puede
    \item nosotros/as - podemos
    \item vosotros/as - podéis
    \item ustedes - pueden
    \item ellos/ellas - pueden
\end{itemize}

\subsection*{Pretérito Indefinido}
\begin{itemize}[leftmargin=*]
    \item yo - pude
    \item tú - pudiste
    \item él/ella/usted - pudo
    \item nosotros/as - pudimos
    \item vosotros/as - pudisteis
    \item ustedes - pudieron
    \item ellos/ellas - pudieron
\end{itemize}

\subsection*{Pretérito Imperfecto}
\begin{itemize}[leftmargin=*]
    \item yo - podía
    \item tú - podías
    \item él/ella/usted - podía
    \item nosotros/as - podíamos
    \item vosotros/as - podíais
    \item ustedes - podían
    \item ellos/ellas - podían
\end{itemize}

\columnbreak

\subsection*{Futuro Simple}
\begin{itemize}[leftmargin=*]
    \item yo - podré
    \item tú - podrás
    \item él/ella/usted - podrá
    \item nosotros/as - podremos
    \item vosotros/as - podréis
    \item ustedes - podrán
    \item ellos/ellas - podrán
\end{itemize}

\subsection*{Pretérito Perfecto}
\begin{itemize}[leftmargin=*]
    \item yo - he podido
    \item tú - has podido
    \item él/ella/usted - ha podido
    \item nosotros/as - hemos podido
    \item vosotros/as - habéis podido
    \item ustedes - han podido
    \item ellos/ellas - han podido
\end{itemize}

\end{multicols}

\newpage

\section*{Explicación Gramatical}

\textbf{Reflexive Verben und Gerundio}

\textbf{1. Reflexive Verben}

Reflexive Verben sind Verben, bei denen die Handlung auf das Subjekt zurückfällt. Sie werden mit einem Reflexivpronomen verwendet (me, te, se, nos, os, se).

Die Reflexivpronomen sind:
\begin{itemize}[leftmargin=*]
    \item \textbf{me} - mich/mir
    \item \textbf{te} - dich/dir
    \item \textbf{se} - sich
    \item \textbf{nos} - uns
    \item \textbf{os} - euch (nur in Spanien)
    \item \textbf{se} - sich (Plural)
\end{itemize}

Reflexive Verben stehen normalerweise vor dem konjugierten Verb oder werden an den Infinitiv oder das Gerundium angehängt.

\textbf{2. Gerundio (Verlaufsform)}

Das Gerundio entspricht dem deutschen Partizip Präsens (laufend, schwimmend) oder wird für die Verlaufsform verwendet (estoy corriendo = ich laufe gerade).

Die Bildung des Gerundio:
\begin{itemize}[leftmargin=*]
    \item \textbf{-ar Verben:} -ar → -ando (hablar → hablando)
    \item \textbf{-er Verben:} -er → -iendo (correr → corriendo)
    \item \textbf{-ir Verben:} -ir → -iendo (vivir → viviendo)
\end{itemize}

Einige Verben haben unregelmäßige Gerundio-Formen:
\begin{itemize}[leftmargin=*]
    \item \textbf{ir} → yendo
    \item \textbf{leer} → leyendo
    \item \textbf{oír} → oyendo
    \item \textbf{pedir} → pidiendo
    \item \textbf{seguir} → siguiendo
    \item \textbf{dormir} → durmiendo
    \item \textbf{venir} → viniendo
\end{itemize}

\textbf{3. Gerundio als Adverbiale Bestimmung}

Das Gerundio kann auch verwendet werden, um eine gleichzeitige Handlung oder einen Grund auszudrücken (ähnlich wie im Deutschen: \enquote{laufend}, \enquote{schwimmend}).

\textbf{Ejemplos del texto:}

\textbf{Reflexive Verben:}
\begin{itemize}[leftmargin=*]
    \item \enquote{Me levanto temprano} -- Ich stehe früh auf
    \item \enquote{Me despierto a las seis} -- Ich wache um sechs auf
    \item \enquote{Me ducho} -- Ich dusche (mich)
    \item \enquote{Me visto} -- Ich ziehe mich an
    \item \enquote{Me siento muy bien} -- Ich fühle mich sehr gut
    \item \enquote{Me acostumbré} -- Ich gewöhnte mich daran
    \item \enquote{Me estiro} -- Ich dehne mich
    \item \enquote{Me alimento muy bien} -- Ich ernähre mich sehr gut
    \item \enquote{Me siento satisfecha} -- Ich fühle mich zufrieden
    \item \enquote{Me relajo} -- Ich entspanne mich
    \item \enquote{Me calmo} -- Ich beruhige mich
    \item \enquote{Me pongo metas} -- Ich setze mir Ziele
    \item \enquote{Me premio} -- Ich belohne mich
    \item \enquote{Te levantas} -- Du stehst auf
    \item \enquote{Te sientes estresada} -- Du fühlst dich gestresst
\end{itemize}

\textbf{Gerundio:}
\begin{itemize}[leftmargin=*]
    \item \enquote{He estado haciendo mucho ejercicio} -- Ich habe viel Sport gemacht (Verlaufsform)
    \item \enquote{Voy corriendo al parque} -- Ich gehe laufend zum Park (Gerundio als Adverbiale)
    \item \enquote{Corriendo me siento más enérgica} -- Laufend fühle ich mich energischer (Gerundio als Adverbiale)
    \item \enquote{Haciendo ejercicio regularmente, duermo mejor} -- Indem ich regelmäßig Sport mache, schlafe ich besser (Gerundio als Adverbiale)
    \item \enquote{Terminando el ejercicio, me estiro} -- Nachdem ich den Sport beendet habe, dehne ich mich (Gerundio als Adverbiale)
    \item \enquote{Voy nadando dos veces por semana} -- Ich gehe zweimal pro Woche schwimmen
    \item \enquote{Nadando trabajo todos los músculos} -- Beim Schwimmen trainiere ich alle Muskeln
    \item \enquote{Haciendo yoga, me relajo} -- Beim Yoga mache ich entspanne ich mich
    \item \enquote{Almorzando verduras, me siento satisfecha} -- Indem ich Gemüse zu Mittag esse, fühle ich mich zufrieden
    \item \enquote{Cenando temprano, duermo mejor} -- Indem ich früh zu Abend esse, schlafe ich besser
    \item \enquote{Caminando media hora, ya estás haciendo algo bueno} -- Indem du eine halbe Stunde gehst, tust du schon etwas Gutes
    \item \enquote{Bebiendo agua, tu cuerpo funciona mejor} -- Indem du Wasser trinkst, funktioniert dein Körper besser
    \item \enquote{Comiendo alimentos saludables, te sentirás mejor} -- Indem du gesunde Lebensmittel isst, wirst du dich besser fühlen
    \item \enquote{Haciendo ejercicios en casa, te mantienes activo} -- Indem du zu Hause Übungen machst, bleibst du aktiv
    \item \enquote{Acostándote temprano, descansas más} -- Indem du früh ins Bett gehst, ruhst du dich mehr aus
    \item \enquote{Evitando las pantallas, duermes mejor} -- Indem du Bildschirme vermeidest, schläfst du besser
    \item \enquote{Respirando profundamente, me calmo} -- Indem ich tief atme, beruhige ich mich
    \item \enquote{Escuchando música, me siento mejor} -- Indem ich Musik höre, fühle ich mich besser
    \item \enquote{Poniéndome metas, me siento motivada} -- Indem ich mir Ziele setze, fühle ich mich motiviert
    \item \enquote{Premiándome, me mantengo entusiasmada} -- Indem ich mich belohne, bleibe ich begeistert
    \item \enquote{Haciendo algo ligero, no pierdes el hábito} -- Indem du etwas Leichtes machst, verlierst du die Gewohnheit nicht
    \item \enquote{Empezando poco a poco, puedes lograr grandes cosas} -- Indem du Schritt für Schritt anfängst, kannst du große Dinge erreichen
\end{itemize}

\textbf{Wichtige Hinweise:}
\begin{itemize}[leftmargin=*]
    \item Reflexive Verben werden oft für alltägliche Handlungen verwendet (sich waschen, sich anziehen, sich setzen)
    \item Das Gerundio kann sowohl für die Verlaufsform (estoy haciendo) als auch als Adverbiale Bestimmung verwendet werden
    \item Wenn ein Reflexivpronomen an ein Gerundium angehängt wird, kann ein Akzent hinzugefügt werden müssen: \enquote{estirándome} statt \enquote{estirando me}
    \item In südamerikanischem Spanisch wird \enquote{os} nicht verwendet, stattdessen \enquote{se} für die 2. Person Plural
\end{itemize}

\end{document}
